\section{Memory Hierarchy: Caches, Main Memory, Virtual Memory}\label{ch:mem}

\subsection{The framework: AMAT and the four levers}
\[
\AMAT = t_{hit}+\text{miss rate}\times\text{miss penalty},\qquad
\CPI_{total}=\CPI_{base}+\frac{\text{mem.\ refs}}{\text{instr.}}\times\text{miss rate}\times\text{miss penalty}\bs
\]
(The slides use exactly this CPI form in the memory-width example: $\CPI=2+1.2\times3\%\times64=4.3$.)

\begin{center}
\small
\begin{tabular}{L{3.8cm}L{11.6cm}}
\toprule
\tblhead{\textbf{Lever} & \textbf{Techniques (slides)}}\\
\midrule
Reduce \textbf{miss penalty} & Multilevel caches (the most important), critical word first \& early restart, read-miss priority over writes, merging write buffer, victim caches \\
Reduce \textbf{miss rate} & Larger blocks, larger caches, higher associativity, way prediction, pseudo-associative caches, compiler optimisations \\
Reduce penalty/rate via \textbf{parallelism} & Non-blocking caches, hardware prefetch, compiler-controlled prefetch \\
Reduce \textbf{hit time} & Small \& simple caches, avoid address translation in indexing, pipelined cache access, trace caches \\
\bottomrule
\end{tabular}
\end{center}
\textbf{Why hit time matters most:} it affects the clock cycle, so a fast hit helps \emph{everything} and its importance is ``multiplied'' beyond the AMAT formula; also miss penalty (in cycles) depends on clock cycle time.

\subsection{Step 1: diagnose the miss type (the ``3 C's'' + coherence)}
\small
\begin{longtable}{L{2.6cm}L{4.4cm}L{4.4cm}L{4.0cm}}
\toprule
\tblhead{\textbf{Miss type} & \textbf{Definition} & \textbf{Reduced by} & \textbf{Does NOT fix}}\\
\midrule
\endhead
Compulsory (cold) & First access to a block; measured by an infinite cache; independent of cache size & Larger blocks (spatial locality), prefetching & Bigger cache or more associativity \\
Capacity & Cache cannot hold all blocks needed; measured by fully associative cache of same size; falls as cache grows & Larger cache; blocking/tiling (shrinks working set); better temporal locality & Associativity \\
Conflict (collision) & Too many blocks map to the same set; vanishes with full associativity; falls as associativity rises & Higher associativity, victim cache, pseudo-associative, way-predicted set-assoc., code/data layout & Compulsory/capacity \\
Coherence (multiprocessor) & Block invalidated by another processor's write. \textbf{True sharing:} the same word is written by another processor. \textbf{False sharing:} different words that happen to be in the same cache line & Data layout (separate independently-updated data into different lines)\bs; smaller blocks\bs & Bigger caches \\
\bottomrule
\end{longtable}
\normalsize
\textbf{Slide examples:} 32\,B lines, array of 32 doubles: first touch of \texttt{array[0]} is a cold miss; prefetching \texttt{array[1..3]} turns later accesses into hits. In a direct-mapped cache, \texttt{array[0]}$\to$\texttt{array[16]}$\to$\texttt{array[0]} conflicts even though the cache is not full. In a multiprocessor, P2 writing \texttt{array[1]} invalidates P1's line holding \texttt{array[0]}: false-sharing miss.
\src{ca05 ``Classify Cache Misses --- 3 C's''; Cache-Misses (``Taxonomy of Cache Misses'')}

\subsection{Reducing miss penalty}
\subsubsection{Multilevel caches (``probably the best miss-penalty reduction'')}
\[
\AMAT=t_{hit,L1}+m_{L1}\big(t_{hit,L2}+m_{L2}^{local}\times P_{L2}\big),\quad
\text{global miss rate}=m_{L1}\times m_{L2}^{local}
\]
\begin{itemize}
\item \textbf{Local} miss rate $=$ misses in this cache / accesses \emph{to this cache}. \textbf{Global} $=$ misses / accesses generated \emph{by the CPU}. \textbf{Global miss rate is what matters} when judging L2/L3.
\item L1 capacity \emph{and} conflict misses are served by L2, avoiding main memory. A big L2 ($\gg$L1) gives a global miss rate close to a single-level cache of that size; the slide's conclusion is that L2 \emph{cache size} is what matters.
\item \textbf{L2 associativity:} direct-mapped hit 10\,cc, local miss 25\%; 2-way 10.1$\to$10 or 11\,cc, local miss 20\%; $P_{L2}=100$\,cc: DM $=10+0.25\times100=35$; 2-way $=10+0.20\times100=30$ (31 if the hit costs 11). \textbf{Higher associativity pays off at L2} because the extra hit time is tiny compared to the miss penalty.
\end{itemize}
\begin{ex}
1000 refs, 40 L1 misses, 20 L2 misses. $m_{L1}=4\%$; local $m_{L2}=20/40=50\%$; global $=2\%$. With $t_{L1}=1$, $t_{L2}=10$, $P_{L2}=100$, 1.5 refs/instr: $\AMAT=1+0.04(10+0.5\times100)=3.4$\,cc; stall cycles/instr $=(3.4-1)\times1.5=3.6$.
\end{ex}
\begin{pick}
Almost always. Essential when the processor is fast and memory is far (servers, desktops: L1+L2+L3). PMDs use two levels (64\,KB L1 + 256\,KB L2) to save energy/area. Hardware complexity 2 (costly if block sizes differ between levels).
\end{pick}

\subsubsection{Critical word first and early restart}
Do not wait for the whole block. \textbf{Critical word first:} fetch the missed word first and send it to the CPU, then fill the rest of the block. \textbf{Early restart:} fetch in normal order but restart the CPU the moment the requested word arrives. \textbf{Value depends on block size and spatial locality:} useful mainly for \emph{large blocks}; if the next access is the next sequential word in the block, the benefit shrinks.

\subsubsection{Read-miss priority over writes; merging write buffer}
\begin{itemize}
\item \textbf{Write-through with a write buffer:} a read miss may need data still sitting in the buffer (RAW through memory). Rather than waiting for the buffer to empty (increased read-miss penalty; one old MIPS 1000 design with a 4-word buffer saw about 50\% worse), \textbf{check the buffer for an address match and let the read go ahead if there is none}.
\item \textbf{Write-back:} on a read miss that replaces a dirty block, copy the dirty block to a buffer, do the read first, then write the dirty block: the CPU restarts sooner.
\item \textbf{Merging write buffer:} an entry holds several words but a write is usually one word. Merge a new write into an existing valid entry with a matching address: fewer full-buffer stalls and faster multiword writes. Mainly used with write-through caches.
\end{itemize}

\subsubsection{Victim cache}
A small \textbf{fully-associative} cache between the cache and its refill path holding only blocks \emph{just evicted}. On a miss, check it before going to the next level; if present, \textbf{swap} victim and cache block. No penalty increase because both are probed together.
\begin{itemize}
\item Jouppi: removes \textbf{20--95\% of the misses} of a \textbf{4\,KB direct-mapped} cache with 1--5 victim entries. AMD Athlon has 8 entries.
\end{itemize}
\begin{pick}
A direct-mapped (small, fast, simple) L1 suffering \emph{conflict} misses --- you keep the hit time of direct-mapped and recover much of 2-way's miss rate. Hardware complexity 2.
\end{pick}

\subsection{Reducing miss rate}
\subsubsection{Larger block size}
Reduces compulsory misses (spatial locality) but \textbf{raises miss penalty} and, in small caches, conflict and capacity misses. Pick the block that minimises \textbf{AMAT}, not miss rate.

\begin{center}
\small
\begin{tabular}{C{1.9cm}C{2.0cm}C{1.9cm}C{1.9cm}C{1.9cm}C{1.9cm}}
\toprule
\tblhead{\textbf{Block (B)} & \textbf{Penalty (cc)} & \textbf{4K} & \textbf{16K} & \textbf{64K} & \textbf{256K}}\\
\midrule
16 & 82 & 8.027 & 4.231 & 2.673 & 1.894 \\
32 & 84 & \textbf{7.082} & 3.411 & 2.134 & 1.588 \\
64 & 88 & 7.160 & \textbf{3.323} & \textbf{1.933} & \textbf{1.449} \\
128 & 96 & 8.469 & 3.659 & 1.979 & 1.470 \\
256 & 112 & 11.651 & 4.685 & 2.288 & 1.549 \\
\bottomrule
\end{tabular}
\end{center}
AMAT (cc) vs.\ block size; memory model: 80\,cc overhead + 2\,cc per 16\,B. \textbf{32--64\,B blocks dominate}; a 256\,B block only wins once the cache is 256\,KB (miss rate vs.\ block table: 4\,KB cache 256\,B block 9.51\% vs.\ 32\,B 7.24\%). Pentium 4's L2 uses 128\,B blocks.

\begin{dec}[Block size]
Small L1 $\rightarrow$ 32--64\,B. Large L2/L3 $\rightarrow$ 64--128\,B. Higher memory latency or bandwidth (overhead dominates the penalty) $\rightarrow$ somewhat larger blocks\bs. Small cache with 128/256\,B blocks $\rightarrow$ wrong (conflict and capacity misses rise).
\end{dec}

\subsubsection{Larger caches and higher associativity}
\begin{itemize}
\item \textbf{Larger cache:} helps capacity and conflict misses; costs hit time, area, power; popular for off-chip/L2.
\item \textbf{Associativity:} 8-way is nearly as good as fully associative. \textbf{2:1 rule:} 2-way of size $N/2$ $\approx$ direct-mapped of size $N$ (for caches $<$128\,KB). Cost: longer hit time $\Rightarrow$ clock cycle $\times1.10$ (2-way), $\times1.12$ (4-way), $\times1.14$ (8-way) relative to 1-way.
\end{itemize}
\begin{center}
\small
\begin{tabular}{C{2.4cm}C{2.2cm}C{2.2cm}C{2.2cm}C{2.2cm}}
\toprule
\tblhead{\textbf{Size (KiB)} & \textbf{1-way} & \textbf{2-way} & \textbf{4-way} & \textbf{8-way}}\\
\midrule
4 & 3.44 & 3.25 & \textbf{3.22} & 3.28 \\
8 & 2.69 & 2.58 & \textbf{2.55} & 2.62 \\
16 & \textbf{2.23} & 2.40 & 2.46 & 2.53 \\
32 & \textbf{2.06} & 2.30 & 2.37 & 2.45 \\
64 & \textbf{1.92} & 2.14 & 2.18 & 2.25 \\
128 & \textbf{1.52} & 1.84 & 1.92 & 2.00 \\
256 & \textbf{1.32} & 1.66 & 1.74 & 1.82 \\
512 & \textbf{1.20} & 1.55 & 1.59 & 1.66 \\
\bottomrule
\end{tabular}
\end{center}
AMAT with the clock-stretch penalty included (the bold cell is the best per row; in the slide the \emph{worse-than-left} cells are highlighted). \textbf{Reading:} for tiny (4--8\,KB) caches 2--4 way wins; \textbf{from 16\,KB upward direct-mapped has the lowest AMAT} once the slower clock is charged. This is the quantitative reason an L1 that sets the clock is small and low-associativity, while an L2 (off the critical path) can be highly associative.

\subsubsection{Way prediction}
Keep extra bits per set predicting which way will hit next; the multiplexer is set early and \textbf{only one tag comparison is done} first. A wrong prediction costs extra cycle(s) to check the remaining ways. Alpha 21264 uses it in its 2-way I-cache with $>85\%$ accuracy (SPEC95). Predictor styles: last-used way, MRU, history table. Metrics: accuracy, misprediction rate, energy reduction.
\begin{pick}
Set-associative caches where \textbf{energy and hit time} matter: PMDs, instruction caches (sequential access makes prediction accurate), cores where a full parallel tag compare is too slow/hot. It keeps most of the associativity benefit while activating one way. Cost: predictor hardware, an extra access on a misprediction.
\end{pick}

\subsubsection{Pseudo-associative (column-associative) cache}
Start as a \textbf{direct-mapped} cache; on a miss, probe a second location (invert the high-order index bit) \textbf{sequentially}, not in parallel. Gives fast hit (primary) and slow hit (alternate); swap blocks so the frequent one sits in the fast slot. \textbf{Problem:} a variable 1-or-2-cycle hit is hard for a CPU pipeline $\Rightarrow$ better for \textbf{caches not tied directly to the processor (L2)}; used in the L2 of the MIPS R10000 (and similar in UltraSPARC).
\begin{ex}[AMAT, slide values]
2\,KB: 1-way 5.90, 2-way 4.90, pseudo 4.95. 128\,KB: 1-way 1.50, 2-way 1.45, pseudo 1.371. Pseudo-associativity gets (nearly) 2-way miss rate at direct-mapped hit speed.
\end{ex}


\subsubsection{Compiler optimisations (zero hardware cost)}
\begin{itemize}
\item \textbf{Code:} reorder procedures (profile-guided) to cut conflicts; McFarling: $-50\%$ misses for a 2\,KB direct-mapped I-cache with 4\,B blocks, $-75\%$ for 8\,KB; optimised programs on a direct-mapped cache miss less than unoptimised on 8-way of the same size.
\item \textbf{Merging arrays:} combine parallel arrays (\texttt{val[]}, \texttt{key[]}) into an array of structs $\Rightarrow$ better spatial locality, fewer conflicts.
\item \textbf{Loop interchange:} change loop order so the inner loop walks memory sequentially (row-major). $1024\times1024$ ints, 64\,B lines (16 ints): \textbf{3{,}145{,}728 $\rightarrow$ 196{,}608 misses for $C=A+B$ ($-93.75\%$)}.
\item \textbf{Loop fusion:} merge loops over the same data so a value is reused while in cache (temporal locality). Example: A, B, C larger than the cache: $262{,}144 \rightarrow 196{,}608$ line transfers ($-25\%$).
\item \textbf{Blocking (tiling):} operate on $B\times B$ sub-matrices that fit in cache. Worst-case capacity misses $2N^3+N^2 \rightarrow 2N^3/B+N^2$; also helps register allocation. Slide example: $N=1024$, 32\,KB cache, $B=32$ (three tiles $=3\times4$\,KB$=12$\,KB fits): $\approx1.07$\,billion misses $\rightarrow\approx2.1$ million (simplified model).
\end{itemize}
\begin{pick}
Array/loop-dominated code with large data sets (scientific, image, matrix); embedded systems where adding hardware is impossible. \emph{Matrix multiply with arrays accessed both by rows and columns: interchange does not help $\Rightarrow$ use blocking.}
\end{pick}
\begin{avoid}
Pointer-chasing, irregular data structures (no static pattern); loops where interchange/fusion is illegal because of dependences.
\end{avoid}
\src{ca05 ``Reducing Miss Rate'' through ``Summary of Compiler Optimizations'', ``Way Prediction'', ``Pseudo-Associative Caches''; CachePerformance2 (way prediction, pseudo-associative, loop interchange/fusion/blocking numerics)}

\subsection{Reducing penalty/rate through parallelism}
\subsubsection{Non-blocking (lockup-free) cache}
Cache continues to supply hits during a miss. \textbf{Hit-under-miss} (one outstanding miss), \textbf{hit-under-multiple-miss / miss-under-miss} (overlap several misses; Pentium Pro allows 4). \textbf{Requires an out-of-order CPU} (scoreboard/Tomasulo) and \textbf{multiple memory banks}; greatly increases the cache controller complexity; must detect repeat misses to a block already being fetched.
\begin{itemize}
\item OoO CPUs can hide L1 misses that hit in L2, but \textbf{not a significant fraction of L2 misses}.
\item Example (8\,KB, FP): miss rate$\times$penalty: DM $11.4\%\times16=1.84$; 2-way $10.7\%\times16=1.71$. For integer: DM $7.4\%\times16=1.18$, 2-way $6.0\%\times16=0.96$. Hit-under-miss does \textbf{not increase hit time}, unlike associativity.
\item Hard to evaluate: effective miss penalty is the \emph{non-overlapped} time the CPU stalls.
\end{itemize}
\begin{pick}
Any out-of-order superscalar, especially memory-intensive FP code; pair with banked/interleaved memory. Useless on a strict in-order blocking CPU.
\end{pick}

\subsubsection{Hardware prefetching (stream buffers)}
Fetch what you expect to need before it is requested, outside the cache. Alpha AXP 21064: on an I-miss fetch 2 blocks (target to I-cache, next to the \textbf{instruction stream buffer}); 1/4/16-block ISB catches 15--25\%/50\%/72\% of misses. Data: one data stream buffer removed 25\% of a 4\,KB cache's misses, 4 streams 43\%; 8 streams 50--70\% (for two 64\,KB 4-way caches, scientific programs). \textbf{Needs spare memory bandwidth} (otherwise it hurts).
\begin{ex}[Slide]
Hit time 2\,cc, prefetch-hit time 1\,cc, prefetch-hit rate 25\%, miss rate 1.10\% (8\,KB), miss penalty 50\,cc:\\
$\AMAT=2+0.011\times0.25\times1+0.011\times0.75\times50=2.41525$\,cc. A cache without prefetching needs a miss rate of only 0.83\% (between 8\,KB at 1.10\% and 16\,KB at 0.64\%) to match.
\end{ex}

\subsubsection{Compiler-controlled prefetching}
Special prefetch instructions (register prefetch: HP PA-RISC; cache prefetch: MIPS IV, PowerPC, SPARC v9), non-faulting (a form of speculative execution), best for loops (e.g.\ \texttt{prefetch(b[j+7][0])}). Cost: issuing prefetches takes instruction slots, so benefit must exceed issue cost; needs a non-blocking cache. Works for instructions too.

\begin{dec}[Hardware vs.\ compiler prefetch]
Hardware (stream buffers): sequential/streaming access, no compiler support, spare bandwidth. Compiler: regular strided loops where you know the distance; embedded systems with controllable code. Skip prefetching on bandwidth-starved or energy-limited designs.
\end{dec}
\src{ca05 ``Reducing Cache Miss Penalty or Miss Rate Via Parallelism'' through ``Compiler-Controlled Prefetching''}

\subsection{Reducing hit time}
\begin{itemize}
\item \textbf{Small and simple caches:} the slow part of a hit is index $\rightarrow$ read tag $\rightarrow$ compare. Keep L1 small enough to be on-chip; direct-mapped is simpler/faster (more misses but faster hit); for L2 keep tags on-chip and data off-chip. Used when the cache sets the clock cycle.
\item \textbf{Avoid address translation during indexing.} \emph{Physically addressed} caches translate with the TLB first; a \emph{virtually addressed} cache indexes in parallel with translation. Problems of pure virtual caches:
 \begin{itemize}
 \item \textbf{Protection} bits must be copied on miss and checked every access.
 \item \textbf{Process switch} $\Rightarrow$ same virtual address $\neq$ same physical data $\Rightarrow$ flush (huge task-switch cost and compulsory misses); fix: \textbf{PID in the tag} (PIDs save 0.6--4.3\% miss rate over purging; cost 0.3--0.5\% for uniprocess).
 \item \textbf{Synonyms/aliases:} two virtual addresses for one physical address $\Rightarrow$ two copies; fix: hardware anti-aliasing (check PA on miss, invalidate), or \textbf{page colouring} (aliases share low-order bits, e.g.\ 18 bits).
 \item \textbf{I/O} uses physical addresses; needs mapping to VA.
 \end{itemize}
 Compromise from the slides: overlap cache access with VA translation, requiring the cache index to remain invariant across translation (trivial if the cache is small); used in Alpha 21164 and UltraSPARC III.
\item \textbf{Pipelined cache access / pipelined writes:} write hits normally need tag check before data write; pipeline into two stages so the \emph{current write's tag check overlaps the previous write's data update}; looks like one write per cycle, short cycle time. Works if the CPU does not immediately depend on the written data (watch read/write ordering). Reads are unaffected (tag check in parallel with data read).
\item \textbf{Trace cache:} stores dynamic instruction \emph{traces} from branch entry to exit, avoiding wasted space in conventional blocks (a taken branch every 5--10 instructions means much of a 16-instruction/64\,B block is unused). Costs: complicated address mapping, duplicate storage of the same instructions in different traces. Used in Pentium 4 (NetBurst).
\end{itemize}
\begin{pick}
Small simple L1 / pipelined writes: every high-clock design. Virtual/overlapped translation: tight L1 designs with a small cache. Trace cache: wide-issue front ends that need fetch bandwidth across taken branches (not energy-constrained designs).
\end{pick}
\src{ca05 ``Reducing Hit Time'' through ``Trace Caches (Cont.)''}

\subsection{Cache optimisation summary table (from the slides)}
\small
\begin{longtable}{L{4.6cm}C{1.2cm}C{1.2cm}C{1.2cm}C{1.8cm}L{5.2cm}}
\toprule
\tblhead{\textbf{Technique} & \textbf{Miss pen.} & \textbf{Miss rate} & \textbf{Hit time} & \textbf{HW compl.} & \textbf{Comment}}\\
\midrule
\endhead
Multilevel caches & $+$ & & & 2 & Costly; widely used \\
Critical word first / early restart & $+$ & & & 2 & Widely used \\
Read priority over writes & $+$ & & & 1 & Trivial for uniprocessor \\
Merging write buffer & $+$ & & & 1 & Used with write-through \\
Victim caches & $+$ & $+$ & & 2 & AMD Athlon: 8 entries \\
Larger block size & $-$ & $+$ & & 0 & Trivial; Pentium 4 L2 128\,B \\
Larger cache size & & $+$ & $-$ & 1 & Widely used, esp.\ L2 \\
Higher associativity & & $+$ & $-$ & 1 & Widely used \\
Way-predicting caches & & & $+$ & 2 & UltraSPARC III I-cache \\
Pseudo-associative & & $+$ & & 2 & MIPS R10000 L2 \\
Compiler techniques (miss rate) & & $+$ & & 0 & Software is the challenge \\
Non-blocking caches & $+$ & & & 3 & All out-of-order CPUs \\
HW prefetch (instr./data) & $+$ & $+$ & & 2/3 & UltraSPARC III \\
Compiler-controlled prefetch & $+$ & $+$ & & 3 & Needs non-blocking cache \\
Small and simple caches & & & $+$ ($-$ miss rate) & 0 & Trivial; widely used \\
Avoid translation in indexing & & & $+$ & 2 & Trivial if small; Alpha 21164 \\
Pipelined cache access & & & $+$ & 1 & Widely used \\
Trace cache & & & $+$ & 3 & Pentium 4 \\
\bottomrule
\end{longtable}
\normalsize
($+$ improves, $-$ hurts; HW complexity 0 = none, 3 = highest.) Note the \emph{conflicts of interest}: larger blocks hurt penalty; larger size/associativity hurt hit time; small/simple caches hurt miss rate.

\subsection{Main memory}
\subsubsection{DRAM vs.\ SRAM and key metrics}
\begin{itemize}
\item \textbf{Main memory is DRAM:} dynamic, needs periodic refresh (about 8\,ms, $\sim$1\% of time); address split into RAS and CAS halves (2-D array). \textbf{Caches are SRAM:} no refresh; 6 transistors/bit vs.\ 1. DRAM/SRAM size ratio 4--8; SRAM/DRAM cost and cycle-time ratio 8--16.
\item \textbf{Latency:} access time (request $\rightarrow$ word arrives) and cycle time (minimum between requests, $>$ access time) $\Rightarrow$ miss penalty. \textbf{Bandwidth:} bytes per unit time $\Rightarrow$ time to transfer the block (matters for I/O and large-block L2 misses).
\item DDR trend (slide tables): total best-case access time improved only from 42\,ns (DDR1, 2000) to 26\,ns (DDR4, 2016), whereas DIMM bandwidth rose from 2128 (PC2100) to 21{,}300 MiB/s (PC21300) --- again \textbf{bandwidth far outgrows latency}.
\end{itemize}

\subsubsection{Three ways to raise memory bandwidth}
\begin{center}
\small
\begin{tabular}{L{2.4cm}L{4.6cm}L{4.2cm}L{4.4cm}}
\toprule
\tblhead{\textbf{Organisation} & \textbf{What it is} & \textbf{Pros} & \textbf{Cons}}\\
\midrule
Simple & CPU, cache, bus, memory all one word wide & Cheapest & Block transfer is serial: penalty $\propto$ words per block \\
Wide & Memory/bus (and cache side) $N$ words wide (Alpha 64/256\,bit, UltraSPARC 512\,bit) & Bandwidth $\times N$; penalty drops accordingly & Wider bus cost; MUX between cache and CPU can be on critical path (put MUX between L1 and L2); ECC harder (partial update needs read of the rest); expansion in larger increments \\
Interleaved & Same width bus, $N$ memory banks; bank $=$ address mod $\#$banks & Overlaps bank accesses, shares bus; great for \textbf{sequential} access (cache blocks); cheap & Needs $\ge$ (access time in cc) banks to deliver one word/cycle; many small chips costly; hard to expand \\
Independent banks & Multiple controllers & Multiprocessor/I/O, non-blocking caches (hit under $n$ misses) & Controller complexity \\
\bottomrule
\end{tabular}
\end{center}

\begin{ex}[Slide: what interleaving/width buys]
Base CPI 2, 1.2 mem.\ refs/instr, miss penalty = 4 (address) + 56 (access) + 4 per word. Miss rates: 3\%/2\%/1.2\% for 1/2/4-word blocks.
\begin{center}
\small
\begin{tabular}{L{6.2cm}C{2.3cm}C{2.3cm}C{2.3cm}}
\toprule
\tblhead{\textbf{Configuration} & \textbf{1-word} & \textbf{2-word} & \textbf{4-word}}\\
\midrule
64-bit bus \& memory, no interleaving & 4.30 & 5.07 & 5.69 \\
64-bit bus \& memory, interleaved & --- & 3.63 & \textbf{3.09} \\
128-bit bus \& memory, no interleaving & --- & \textbf{3.54} & 3.84 \\
\bottomrule
\end{tabular}
\end{center}
\textit{Note:} the slide labels the narrow bus ``64-bit'' and the wide bus ``128-bit''; its arithmetic treats narrow as one word per transfer and wide as two words per transfer. \textbf{Reading:} wide/interleaved beat simple by a large margin; with larger blocks \textbf{interleaving} wins (it overlaps accesses on a narrow bus), with 2-word blocks a \textbf{128-bit} bus is marginally better. If the bus cannot be widened (pins/cost), choose interleaving.
\end{ex}
\src{ca05 ``Main Memory'' through ``Independent Memory Banks''}

\subsubsection{Memory hierarchy by device class (slide figures)}
\textbf{PMD:} L1 64\,KB (1\,ns), L2 256\,KB (5--10\,ns), no L3, 1--2\,GB DRAM (50--100\,ns), flash 4--64\,GB. \textbf{Laptop/desktop:} adds L3 (4--8 / 8--32\,MB, 10--20\,ns), 4--64\,GB DRAM, flash. \textbf{Server:} L3 16--64\,MB, 32--256\,GB DRAM, flash 1--16\,TB and disk 16--64\,TB (5--10\,ms).

\subsection{Virtual memory}
\begin{itemize}
\item \textbf{Why:} many processes each with their own address space share a smaller physical memory; need \emph{protection} (a process only touches its own blocks); faster program start (not everything resident); frees the programmer from overlays.
\item \textbf{Pages vs.\ segments:} pages $=$ fixed-size blocks (4--8\,KB), one address (page number $+$ offset), \textbf{no external fragmentation}, allocation anywhere (non-contiguous), no compaction; segments $=$ variable size (1\,B up to $2^{16}$--$2^{32}$\,B), address needs 2 words (segment number + offset). \textbf{Paging is the default choice}.
\item \textbf{Translation:} virtual page number indexes the \textbf{page table} (valid bit, physical page number); valid $=0\Rightarrow$ \textbf{page fault} (OS loads from disk). The page table is in memory, so every access would cost two memory accesses $\Rightarrow$ \textbf{TLB}, a special cache of translations.
\item \textbf{TLB:} usually small ($\le128$--256 entries); high-end: fully associative; mid-range: small $n$-way set associative. Fields: valid, dirty, tag, physical page number. Physical address then drives cache tag compare (TLB $\rightarrow$ cache interaction).
\item Example split: 32-bit VA $=$ 20-bit VPN $+$ 12-bit offset (4\,KB pages).
\end{itemize}
\src{ca05 ``Virtual Memory'' through ``Techniques for Fast Address Translation''}

\subsection{Scenario $\rightarrow$ memory design, quick reference}
\small
\begin{longtable}{L{5.2cm}L{4.1cm}L{6.3cm}}
\toprule
\tblhead{\textbf{Symptom / constraint} & \textbf{Choose} & \textbf{Reason}}\\
\midrule
\endhead
Small direct-mapped L1, ping-pong conflicts & Victim cache; or pseudo-assoc./2-way; or array merging/padding & Conflict misses; victim cache keeps DM hit time. \\
Streaming/sequential large arrays & Stream-buffer prefetch, moderate block (64\,B), interleaved memory, loop interchange & Compulsory/spatial misses; bandwidth is plentiful. \\
Large matrix multiply & Blocking (tiling) with 3 tiles $\le$ cache & Capacity misses; interchange cannot help both row/column accesses. \\
OoO core with many outstanding loads & Non-blocking cache + banked memory & Overlap misses; otherwise OoO core is wasted. \\
L1 hit time limits clock & Small, simple (DM) L1, pipelined writes, VA/translation overlap & Hit time is multiplied across every instruction. \\
PMD / battery & Two cache levels, way prediction, moderate sizes, no trace cache, selective prefetch & Energy and area; way prediction cuts tag-compare energy. \\
Server with large working set & L3 (16--64\,MB), high-associativity L2, interleaved/banked memory & Global miss rate matters; capacity misses dominate. \\
Write-through L1 & Merging write buffer + read priority over writes & Prevent buffer-full stalls and RAW delays. \\
Large blocks and long block fill & Critical word first \& early restart & CPU resumes when the needed word arrives. \\
Instruction-fetch bound, taken branches every 5--10 instr. & Trace cache or I-stream buffer; BTB & Fetch bandwidth across taken branches. \\
Multicore sharing data & Watch coherence (true/false sharing) misses; separate independently-written data into different lines\bs & Coherence misses are a distinct fourth C. \\
Cache must be tiny but fast & Direct-mapped, small block & Table: DM wins AMAT $\ge$16\,KB; for 4--8\,KB use 2--4 way. \\
\bottomrule
\end{longtable}
\normalsize


